\section{代码审查与常见陷阱}

\subsection{三个常见问题}

\begin{warningbox}{Claude Code 代码的典型问题}
\begin{enumerate}[leftmargin=1.5em]
  \item \textbf{过度工程}——你要一个简单工具函数，它给你搞出完整的抽象层（泛型+接口+工厂模式）。\textbf{解法}：在 CLAUDE.md 写"优先用简单方案"，或在指令里明确"不需要抽象"。
  \item \textbf{幻觉 API}——使用不存在的 API 或过时的语法，尤其是小众库的新版本。\textbf{解法}：跑测试，指令里包含验证标准。
  \item \textbf{改动范围膨胀}——你让它改一个函数，它把整个文件都重构了。\textbf{解法}：明确说"只改 X，不要动其他代码"，或用 \texttt{/freeze} 限制编辑范围。
\end{enumerate}
\end{warningbox}

\subsection{Writer/Reviewer 模式}

官方推荐的高级模式：用两个独立会话分别扮演"写代码"和"审查代码"的角色。两个会话有独立上下文，Reviewer 不受 Writer 思路影响，能发现 Writer 的盲点。

\subsection{不要用 Claude Code 做的事}

\begin{warningbox}{四条红线}
\begin{enumerate}[leftmargin=1.5em]
  \item \textbf{不要让它做你完全不了解的事}——你无法审查你不懂的东西。
  \item \textbf{不要在没有版本控制的环境下用}——没有 Git 就没有回退能力，等于裸奔。
  \item \textbf{不要一口气扔一个巨大的任务}——拆成小步骤，每步确认后再做下一步。
  \item \textbf{不要指望一次就完美}——迭代是正常的，用 \texttt{/rewind} 回退，用精确反馈描述"哪里不对"。
\end{enumerate}
\end{warningbox}

\subsection{本章小结}

信任但要验证。用 CLAUDE.md 约束过度工程，用测试拦截幻觉 API，用明确指令控制改动范围。Writer/Reviewer 双会话模式适合重要功能开发。

%% ====================================================================
